\chapter{文本操作的应用}
\begin{lstlisting}[style=pythonstyle, title={删除文本注释.py}]
import re
import os
import readline  # 支持终端 Tab 补全
import glob      # 用于文件路径匹配

def complete_path(text, state):
    """实现终端文件路径 Tab 补全"""
    results = glob.glob(text + '*') + [None]
    return results[state]

# 注册路径补全
readline.set_completer(complete_path)
readline.parse_and_bind("tab: complete")

def extract_comments(file_content, language):
    """
    根据语言类型提取注释。
    :param file_content: 文件内容
    :param language: 用户选择的语言类型
    """
    if language == "C/C++":
        single_line_pattern = r"(?m)^\s*//.*$|(?<=\s)//.*$"
        multi_line_pattern = r"(?s)/\*.*?\*/"
    elif language == "Python":
        single_line_pattern = r"(?m)(^\s*#.*$|\s+#.*$)"
        multi_line_pattern = r"(?s)('''.*?'''|\"\"\".*?\"\"\")"
    elif language == "LaTeX":
        single_line_pattern = r"(?m)^\s*%.*$|(?<=\s)%.*$"
        multi_line_pattern = None

    # 提取注释
    single_line_comments = re.findall(single_line_pattern, file_content)
    multi_line_comments = re.findall(multi_line_pattern, file_content) if multi_line_pattern else []

    return [c for c in (single_line_comments + multi_line_comments) if c.strip()]

def list_comments_with_index(comments):
    """列出所有检测到的注释"""
    print("检测到以下注释：")
    for idx, comment in enumerate(comments, 1):
        print(f"{idx}. {repr(comment.strip())}")

def get_comments_to_keep(comments):
    """用户选择需要保留的注释"""
    user_input = input("输入需要保留的注释编号（用逗号分隔），直接回车全部删除： ")
    if not user_input.strip():
        return []

    try:
        indices = [int(x) - 1 for x in user_input.split(',')]
        return [comments[i] for i in indices if 0 <= i < len(comments)]
    except ValueError:
        print("输入格式不正确，将默认删除所有注释。")
        return []

def remove_comments(file_content, comments_to_keep, all_comments):
    """删除未保留的注释"""
    for comment in all_comments:
        if comment not in comments_to_keep:
            file_content = file_content.replace(comment, '', 1)
    return file_content

def clean_file(input_file, language):
    """主函数：读取文件，处理注释，生成新文件"""
    with open(input_file, 'r', encoding='utf-8') as f:
        file_content = f.read()

    comments = extract_comments(file_content, language)

    if not comments:
        print("未检测到任何注释。")
        return

    list_comments_with_index(comments)
    comments_to_keep = get_comments_to_keep(comments)

    cleaned_content = remove_comments(file_content, comments_to_keep, comments)

    # 生成新文件名并写入
    base_name, ext = os.path.splitext(input_file)
    new_file_name = f"{base_name}_clean{ext}"
    with open(new_file_name, 'w', encoding='utf-8') as f:
        f.write(cleaned_content)

    print(f"处理完成，新文件已保存为：{new_file_name}")

def main():
    """程序入口，用户选择语言和输入文件"""
    print("请选择要处理的语言类型：")
    print("1. C/C++")
    print("2. Python")
    print("3. LaTeX")

    language_choice = input("请输入对应的数字： ").strip()
    language_map = {"1": "C/C++", "2": "Python", "3": "LaTeX"}

    if language_choice not in language_map:
        print("无效选择，请重新运行程序。")
        return

    language = language_map[language_choice]
    input_file = input(f"请输入要处理的 {language} 文件路径： ")

    if os.path.isfile(input_file):
        clean_file(input_file, language)
    else:
        print("文件不存在，请检查路径是否正确。")

if __name__ == "__main__":
    main()
\end{lstlisting}

\begin{lstlisting}[style=pythonstyle, title={递归删除整个文件夹中注释.py}]
import re
import os
import shutil
import readline  # 支持终端 Tab 补全
import glob      # 用于文件路径匹配

def complete_path(text, state):
    """
    实现终端文件路径Tab补全功能。
    
    Args:
        text (str): 用户当前输入的文本
        state (int): 补全状态索引
        
    Returns:
        str: 匹配的路径补全结果，如果没有更多匹配则返回None
    """
    results = glob.glob(text + '*') + [None]
    return results[state]

# 注册路径补全
readline.set_completer(complete_path)
readline.parse_and_bind("tab: complete")

def get_file_language(file_path):
    """
    根据文件后缀名判断编程语言类型。
    
    Args:
        file_path (str): 文件路径
        
    Returns:
        str: 语言类型字符串，如果不支持则返回None
        
    支持的文件类型：
        - C/C++: .c, .cpp, .cc, .cxx, .h, .hpp, .hxx
        - Python: .py, .pyw
        - LaTeX: .tex, .latex  
        - Bash: .sh, .bash
        - Lua: .lua
    """
    ext = os.path.splitext(file_path)[1].lower()
    
    if ext in ['.c', '.cpp', '.cc', '.cxx', '.h', '.hpp', '.hxx']:
        return 'C/C++'
    elif ext in ['.py', '.pyw']:
        return 'Python'
    elif ext in ['.tex', '.latex']:
        return 'LaTeX'
    elif ext in ['.sh', '.bash']:
        return 'Bash'
    elif ext in ['.lua']:
        return 'Lua'
    else:
        return None

def extract_comments(file_content, language):
    """
    根据编程语言类型提取文件中的注释内容。
    
    Args:
        file_content (str): 文件内容
        language (str): 编程语言类型
        
    Returns:
        list: 提取到的注释列表
        
    支持的注释格式：
        - C/C++: 单行注释(//) 和 多行注释(/* */)
        - Python: 单行注释(#) 和 文档字符串(''' ''' 或 """ """)
        - LaTeX: 单行注释(%)
        - Bash: 单行注释(#) 和 多行注释(<< 'EOF' ... EOF)
        - Lua: 单行注释(--) 和 多行注释(--[[ ]])
    """
    if language == "C/C++":
        # C/C++ 单行注释 // 和多行注释 /* */
        single_line_pattern = r"(?m)^\s*//.*$|(?<=\s)//.*$"
        multi_line_pattern = r"(?s)/\*.*?\*/"
    elif language == "Python":
        # Python 单行注释 # 和文档字符串 ''' """ 
        single_line_pattern = r"(?m)(^\s*#.*$|\s+#.*$)"
        multi_line_pattern = r"(?s)('''.*?'''|\"\"\".*?\"\"\")"
    elif language == "LaTeX":
        # LaTeX 单行注释 %
        single_line_pattern = r"(?m)^\s*%.*$|(?<=\s)%.*$"
        multi_line_pattern = None
    elif language == "Bash":
        # Bash 单行注释 # 和多行注释 << 'EOF' ... EOF 或 << "EOF" ... EOF
        single_line_pattern = r"(?m)(^\s*#.*$|\s+#.*$)"
        multi_line_pattern = r"(?s)<<\s*['\"]?(\w+)['\"]?\s*\n.*?\n\1"
    elif language == "Lua":
        # Lua 单行注释 -- 和多行注释 --[[ ]] 或 --[=[ ]=] 等
        single_line_pattern = r"(?m)^\s*--(?!\[).*$|(?<=\s)--(?!\[).*$"
        multi_line_pattern = r"(?s)--\[=*\[.*?\]=*\]"
    else:
        return []

    # 提取注释
    single_line_comments = re.findall(single_line_pattern, file_content)
    multi_line_comments = re.findall(multi_line_pattern, file_content) if multi_line_pattern else []

    return [c for c in (single_line_comments + multi_line_comments) if c.strip()]

def list_comments_with_index(comments):
    """
    列出检测到的所有注释并显示编号，用于用户选择保留哪些注释。
    
    Args:
        comments (list): 注释列表
        
    说明:
        此函数会在控制台打印所有检测到的注释，每个注释前面都有编号，
        方便用户在交互模式下选择需要保留的注释。
    """
    print("检测到以下注释：")
    for idx, comment in enumerate(comments, 1):
        print(f"{idx}. {repr(comment.strip())}")

def get_comments_to_keep(comments):
    """
    交互式获取用户希望保留的注释。
    
    Args:
        comments (list): 注释列表
        
    Returns:
        list: 用户选择保留的注释列表
        
    说明:
        此函数提示用户输入需要保留的注释编号（用逗号分隔），
        如果用户直接回车则删除所有注释。输入格式错误时默认删除所有注释。
    """
    user_input = input("输入需要保留的注释编号（用逗号分隔），直接回车全部删除： ")
    if not user_input.strip():
        return []

    try:
        indices = [int(x) - 1 for x in user_input.split(',')]
        return [comments[i] for i in indices if 0 <= i < len(comments)]
    except ValueError:
        print("输入格式不正确，将默认删除所有注释。")
        return []

def remove_comments(file_content, comments_to_keep, all_comments):
    """
    从文件内容中移除未被保留的注释。
    
    Args:
        file_content (str): 原始文件内容
        comments_to_keep (list): 用户选择保留的注释列表
        all_comments (list): 所有检测到的注释列表
        
    Returns:
        str: 移除指定注释后的文件内容
        
    说明:
        此函数会遍历所有检测到的注释，如果注释不在保留列表中，
        就从文件内容中移除该注释。使用replace方法进行替换，
        每次只替换一个匹配项以避免误删。
    """
    for comment in all_comments:
        if comment not in comments_to_keep:
            file_content = file_content.replace(comment, '', 1)
    return file_content

def clean_file(input_file, output_file, interactive=True):
    """
    处理单个文件的注释清理。
    
    Args:
        input_file (str): 输入文件的完整路径
        output_file (str): 输出文件的完整路径
        interactive (bool): 是否使用交互模式，默认True
        
    Returns:
        bool: 处理成功返回True，失败返回False
        
    说明:
        1. 根据文件后缀名自动识别语言类型
        2. 读取文件内容并提取注释
        3. 在交互模式下让用户选择保留的注释
        4. 移除未保留的注释并写入输出文件
        5. 如果文件不支持或没有注释，将直接复制原文件
    """
    # 根据文件后缀名判断语言类型
    language = get_file_language(input_file)
    if not language:
        print(f"跳过不支持的文件类型：{input_file}")
        return False

    # 尝试读取文件内容
    try:
        with open(input_file, 'r', encoding='utf-8') as f:
            file_content = f.read()
    except Exception as e:
        print(f"读取文件失败 {input_file}: {e}")
        return False

    # 提取注释
    comments = extract_comments(file_content, language)

    # 如果没有注释，直接复制文件
    if not comments:
        print(f"文件 {input_file} 未检测到任何注释，直接复制。")
        try:
            with open(output_file, 'w', encoding='utf-8') as f:
                f.write(file_content)
            return True
        except Exception as e:
            print(f"写入文件失败 {output_file}: {e}")
            return False

    # 根据交互模式处理注释选择
    if interactive:
        print(f"\n处理文件：{input_file} ({language})")
        list_comments_with_index(comments)
        comments_to_keep = get_comments_to_keep(comments)
    else:
        # 非交互模式下删除所有注释
        comments_to_keep = []

    # 移除注释并生成清理后的内容
    cleaned_content = remove_comments(file_content, comments_to_keep, comments)

    # 尝试写入输出文件
    try:
        with open(output_file, 'w', encoding='utf-8') as f:
            f.write(cleaned_content)
        print(f"处理完成：{input_file} -> {output_file}")
        return True
    except Exception as e:
        print(f"写入文件失败 {output_file}: {e}")
        return False

def process_directory(input_dir, output_dir=None, interactive=True):
    """
    递归处理目录中所有支持的文件，移除注释并输出到新目录。
    
    Args:
        input_dir (str): 要处理的输入目录路径
        output_dir (str, optional): 输出目录路径，默认为输入目录名 + '_clean'
        interactive (bool): 是否使用交互模式，默认True
        
    说明:
        1. 递归遍历输入目录中的所有文件和子目录
        2. 对于支持的文件类型（C/C++, Python, LaTeX, Bash, Lua）进行注释清理
        3. 对于不支持的文件类型，直接复制到输出目录
        4. 保持原有的目录结构和文件层次关系
        5. 在交互模式下，每个文件都会让用户选择保留的注释
        6. 如果输出目录已存在，会询问用户是否删除重建
    """
    # 检查输入目录是否存在
    if not os.path.isdir(input_dir):
        print(f"目录不存在：{input_dir}")
        return
    
    # 设置默认输出目录
    if output_dir is None:
        output_dir = f"{input_dir}_clean"
    
    # 处理已存在的输出目录
    if os.path.exists(output_dir):
        response = input(f"输出目录 {output_dir} 已存在，是否删除重建？(y/n): ").strip().lower()
        if response == 'y':
            shutil.rmtree(output_dir)
        else:
            print("操作已取消。")
            return
    
    # 创建输出目录
    os.makedirs(output_dir, exist_ok=True)
    
    # 初始化计数器
    processed_count = 0  # 已处理的文件数量
    skipped_count = 0    # 跳过的文件数量
    
    # 递归遍历输入目录
    for root, dirs, files in os.walk(input_dir):
        # 计算当前目录相对于输入目录的相对路径
        rel_path = os.path.relpath(root, input_dir)
        
        # 确定当前目录对应的输出目录路径
        if rel_path != '.':
            current_output_dir = os.path.join(output_dir, rel_path)
        else:
            current_output_dir = output_dir
        
        # 创建对应的输出目录（如果不存在）
        os.makedirs(current_output_dir, exist_ok=True)
        
        # 处理当前目录下的所有文件
        for file in files:
            input_file = os.path.join(root, file)
            output_file = os.path.join(current_output_dir, file)
            
            # 检查文件是否为支持的类型
            if get_file_language(input_file):
                # 支持的文件类型，进行注释处理
                if clean_file(input_file, output_file, interactive):
                    processed_count += 1
                else:
                    skipped_count += 1
            else:
                # 不支持的文件类型，直接复制到输出目录
                try:
                    shutil.copy2(input_file, output_file)
                    skipped_count += 1
                except Exception as e:
                    print(f"复制文件失败 {input_file}: {e}")
                    skipped_count += 1
    
    # 打印处理结果统计
    print(f"\n处理完成！")
    print(f"已处理文件：{processed_count} 个")
    print(f"跳过文件：{skipped_count} 个")
    print(f"输出目录：{output_dir}")

def main():
    """
    程序主入口函数，负责用户交互和参数设置。
    
    功能说明：
        1. 显示程序介绍和支持的文件类型
        2. 获取用户输入的目录路径
        3. 询问用户是否使用交互模式
        4. 设置输出目录路径
        5. 调用process_directory函数开始处理
        
    支持的文件类型：
        - C/C++: .c, .cpp, .cc, .cxx, .h, .hpp, .hxx
        - Python: .py, .pyw
        - LaTeX: .tex, .latex
        - Bash: .sh, .bash
        - Lua: .lua
        
    交互模式说明：
        - 交互模式(y)：对每个有注释的文件，都会显示所有注释并让用户选择保留哪些
        - 批处理模式(n)：自动删除所有检测到的注释，无需用户干预
    """
    # 显示程序标题和功能介绍
    print("代码注释清理工具 - 目录批处理版本")
    print("支持的文件类型：C/C++ (.c, .cpp, .h, .hpp等), Python (.py), LaTeX (.tex), Bash (.sh), Lua (.lua)")
    print()
    
    # 获取用户输入的目录路径
    input_dir = input("请输入要处理的目录路径：").strip()
    
    # 验证目录是否存在
    if not os.path.isdir(input_dir):
        print("目录不存在，请检查路径是否正确。")
        return
    
    # 询问是否使用交互模式
    print("\n模式选择：")
    print("交互模式(y)：每个文件都会询问保留哪些注释")
    print("批处理模式(n)：自动删除所有注释")
    interactive_choice = input("是否使用交互模式？(y/n，默认n): ").strip().lower()
    interactive = interactive_choice == 'y'
    
    # 设置输出目录
    default_output = f"{input_dir}_clean"
    print(f"\n输出目录设置（默认：{default_output}）")
    output_dir = input(f"请输入输出目录路径（直接回车使用默认）：").strip()
    if not output_dir:
        output_dir = default_output
    
    # 显示处理参数确认
    print(f"\n=== 处理参数确认 ===")
    print(f"输入目录: {input_dir}")
    print(f"输出目录: {output_dir}")
    print(f"处理模式: {'交互模式' if interactive else '批处理模式'}")
    print("=" * 20)
    
    # 开始处理目录
    process_directory(input_dir, output_dir, interactive)

if __name__ == "__main__":
    main()
\end{lstlisting}
\begin{lstlisting}[style=demostyle,numbers=none]
$> python uncomment_dir.py
代码注释清理工具 - 目录批处理版本
支持的文件类型：C/C++ (.c, .cpp, .h, .hpp等), Python (.py), LaTeX (.tex), Bash (.sh), Lua (.lua)

请输入要处理的目录路径：./old

模式选择：
交互模式(y)：每个文件都会询问保留哪些注释
批处理模式(n)：自动删除所有注释
是否使用交互模式？(y/n，默认n):

输出目录设置（默认：./old_clean）
请输入输出目录路径（直接回车使用默认）：

=== 处理参数确认 ===
输入目录: ./old
输出目录: ./old_clean
处理模式: 批处理模式
====================
输出目录 ./old_clean 已存在，是否删除重建？(y/n): y
处理完成：./old/bash/7z_unrar.sh -> ./old_clean/bash/7z_unrar.sh
处理完成：./old/bash/7z_unzip.sh -> ./old_clean/bash/7z_unzip.sh
处理完成：./old/C/zipfolder.c -> ./old_clean/C/zipfolder.c
处理完成：./old/C/zip_split_folder.c -> ./old_clean/C/zip_split_folder.c
处理完成：./old/python/audio2video.py -> ./old_clean/python/audio2video.py
处理完成：./old/python/chinese_num_covert.py -> ./old_clean/python/chinese_num_covert.py
处理完成：./old/python/git_add_batch.py -> ./old_clean/python/git_add_batch.py
处理完成：./old/python/git_batch_folder.py -> ./old_clean/python/git_batch_folder.py
处理完成：./old/python/merged_avi.py -> ./old_clean/python/merged_avi.py
处理完成：./old/python/merged_flv.py -> ./old_clean/python/merged_flv.py
处理完成：./old/python/merged_mp4.py -> ./old_clean/python/merged_mp4.py
处理完成：./old/python/merged_ts.py -> ./old_clean/python/merged_ts.py
文件 ./old/python/name.py 未检测到任何注释，直接复制。
文件 ./old/python/name_dir.py 未检测到任何注释，直接复制。
处理完成：./old/python/uncomment.py -> ./old_clean/python/uncomment.py
处理完成：./old/python/uncomment_dir.py -> ./old_clean/python/uncomment_dir.py
处理完成：./old/python/unzip.py -> ./old_clean/python/unzip.py
处理完成：./old/python/unzip_file.py -> ./old_clean/python/unzip_file.py
处理完成：./old/python/video_timeline.py -> ./old_clean/python/video_timeline.py
处理完成：./old/python/video_timeline_cur_folder.py -> ./old_clean/python/video_timeline_cur_folder.py
处理完成：./old/python/zip.py -> ./old_clean/python/zip.py
处理完成：./old/python/zip100.py -> ./old_clean/python/zip100.py

处理完成！
已处理文件：22 个
跳过文件：4 个
输出目录：./old_clean
\end{lstlisting}

